% ============================================================================
% Slides TCC — Predição de Direção de Preços de Ações Brasileiras
% Apresentação ~15 min | Beamer | pt-BR
% Compilar: pdflatex apresentacao.tex (2x)
% ============================================================================
\documentclass[10pt,aspectratio=169]{beamer}

\usepackage[utf8]{inputenc}
\usepackage[T1]{fontenc}
\usepackage[brazil]{babel}
\usepackage{booktabs}
\usepackage{graphicx}
\usepackage{tikz}
\usetikzlibrary{arrows.meta,positioning}

\graphicspath{{../../9.baselines/}{../../7.model-evaluation/results/}{../../4.finbert-br/}{../../3.model_traning/}}

% --- Tema e cores ---------------------------------------------------------
\usetheme{Madrid}
\usecolortheme{default}

\definecolor{petroleo}{HTML}{0B3142}   % azul-petróleo escuro
\definecolor{verdeagua}{HTML}{1B7F79}  % teal
\definecolor{coral}{HTML}{E4572E}      % destaque
\definecolor{cinzaclaro}{HTML}{F2F4F5}

\setbeamercolor{structure}{fg=petroleo}
\setbeamercolor{palette primary}{bg=petroleo,fg=white}
\setbeamercolor{palette secondary}{bg=verdeagua,fg=white}
\setbeamercolor{palette tertiary}{bg=petroleo,fg=white}
\setbeamercolor{frametitle}{bg=petroleo,fg=white}
\setbeamercolor{title}{bg=petroleo,fg=white}
\setbeamercolor{alerted text}{fg=coral}
\setbeamercolor{block title}{bg=verdeagua,fg=white}
\setbeamercolor{block body}{bg=cinzaclaro,fg=black}

\setbeamertemplate{navigation symbols}{}
\setbeamertemplate{itemize item}{\color{verdeagua}$\blacktriangleright$}
\setbeamertemplate{itemize subitem}{\color{verdeagua}--}
\setbeamerfont{frametitle}{series=\bfseries}

% Rodapé enxuto: autor | título curto | nº slide
\setbeamertemplate{footline}{%
  \leavevmode%
  \hbox{%
    \begin{beamercolorbox}[wd=.45\paperwidth,ht=2.4ex,dp=1ex,left]{palette primary}%
      \hspace*{2ex}\insertshortauthor
    \end{beamercolorbox}%
    \begin{beamercolorbox}[wd=.45\paperwidth,ht=2.4ex,dp=1ex,center]{palette secondary}%
      Sentimento de Notícias e Predição na B3
    \end{beamercolorbox}%
    \begin{beamercolorbox}[wd=.10\paperwidth,ht=2.4ex,dp=1ex,right]{palette primary}%
      \insertframenumber\,/\,\inserttotalframenumber\hspace*{2ex}
    \end{beamercolorbox}}%
}

\title[Sentimento de Notícias e Predição na B3]{Predição de Direção de Preços de Ações Brasileiras\\com Sentimento de Notícias Financeiras}
\subtitle{Pipeline, Ilusão e Autocorreção Metodológica}
\author[André Takeo L. Fujiwara]{André Takeo Loschner Fujiwara\\[1ex]{\small Orientador: Prof.\ Eric Bacconi Gonçalves}}
\institute[PUC-SP]{Ciência de Dados e Inteligência Artificial --- PUC-SP}
\date{Junho de 2026}

\begin{document}

% ---------------------------------------------------------------- 1. Título
\begin{frame}[plain]
  \titlepage
\end{frame}

% ---------------------------------------------------------------- 2. Roteiro
\begin{frame}{Roteiro}
  \begin{columns}[T]
    \column{0.55\textwidth}
    \begin{enumerate}
      \item Contexto e pergunta de pesquisa
      \item Pipeline: notícias + dados de mercado
      \item Resultado inicial: AUC = 0{,}709
      \item Investigação metodológica (1.435 execuções)
      \item Reversão da conclusão e protocolos propostos
    \end{enumerate}
    \column{0.4\textwidth}
    \begin{block}{Mensagem central}
      Um resultado positivo obtido sob protocolo padrão da literatura
      \alert{não sobreviveu} à avaliação rigorosa --- e documentar essa
      reversão é a contribuição principal.
    \end{block}
  \end{columns}
\end{frame}

% ---------------------------------------------------------------- 3. Contexto
\begin{frame}{Contexto e motivação}
  \begin{itemize}
    \item Notícias financeiras como \textit{features} preditivas: linha ativa em NLP + finanças (Xu \& Cohen, 2018; Araci, 2019)
    \item Resultados positivos concentrados em mercados desenvolvidos (S\&P 500)
    \item \textbf{B3 é estruturalmente diferente:} menor liquidez, poucas fontes jornalísticas especializadas, menos participantes institucionais
    \item NLP financeiro em português é recente: FinBERT-PT-BR disponível só a partir de 2022 (Souza et al., 2023)
  \end{itemize}
  \vspace{1ex}
  \begin{block}{Lacuna}
    Pouca evidência empírica para o mercado brasileiro --- e quase nenhuma
    com avaliação metodologicamente rigorosa.
  \end{block}
\end{frame}

% ---------------------------------------------------------------- 4. Pergunta
\begin{frame}{Pergunta de pesquisa}
  \begin{center}
    \Large\textbf{Notícias financeiras brasileiras melhoram a predição\\de direção de preço (sobe/desce) de ações da B3?}
  \end{center}
  \vspace{2ex}
  \begin{columns}[T]
    \column{0.48\textwidth}
    \begin{block}{(a) Representação textual}
      \textit{Embeddings} genéricos de 1.024 dimensões \textbf{ou}
      sentimento financeiro de 5 dimensões (FinBERT-PT-BR)?
    \end{block}
    \column{0.48\textwidth}
    \begin{block}{(b) Robustez metodológica}
      Resultados de \textit{walk-forward split} único sobrevivem a
      CV \textit{expanding-window}, múltiplas sementes e testes formais?
    \end{block}
  \end{columns}
  \vspace{1.5ex}
  \centering
  Estudo de caso: \textbf{ITUB4, PETR4, VALE3} --- horizontes $h=5$ e $h=21$ dias úteis.
\end{frame}

% ---------------------------------------------------------------- 5. Dados
\begin{frame}{Dados: notícias + mercado}
  \begin{columns}[T]
    \column{0.48\textwidth}
    \textbf{Corpus de notícias (InfoMoney)}
    \begin{table}
      \footnotesize
      \begin{tabular}{lrr}
        \toprule
        \textbf{Ativo} & \textbf{Artigos} & \textbf{Período} \\
        \midrule
        ITUB4 (Itaú)      & 2.572 & 2009--2026 \\
        PETR4 (Petrobras) & 1.775 & 2009--2026 \\
        VALE3 (Vale)      & 1.525 & 2009--2026 \\
        \midrule
        \textbf{Total}    & \textbf{5.872} & \\
        \bottomrule
      \end{tabular}
    \end{table}
    \column{0.48\textwidth}
    \textbf{Dados de mercado (Yahoo Finance)}
    \begin{itemize}
      \item OHLCV diário por ativo
      \item 11 \textit{features} técnicas: retorno, médias móveis (7/21), volatilidade (std 21), \textit{lags}
    \end{itemize}
    \textbf{Alvo binário}
    \begin{itemize}
      \item $1$ se $\text{Close}[t+h] > \text{Close}[t]$
      \item Desbalanceamento: 59/41 ``Sobe''/``Desce'' \emph{(treino)} --- compensado por \alert{peso} (\texttt{pos\_weight}/\texttt{scale\_pos\_weight}/peso de amostra conforme o modelo; TCN via \texttt{BCELoss}), \alert{não por reamostragem} --- preserva ordem temporal e \textit{base rate}
    \end{itemize}
  \end{columns}
\end{frame}

% ---------------------------------------------------------------- 6. Pipeline
\begin{frame}{Pipeline}
  \centering
  \resizebox{\textwidth}{!}{%
  \begin{tikzpicture}[
    node distance=4mm and 6mm,
    caixa/.style={rectangle, rounded corners=2pt, draw=petroleo, thick,
                  fill=cinzaclaro, text width=26mm, align=center,
                  minimum height=11mm, font=\scriptsize},
    seta/.style={-{Stealth[length=2.5mm]}, thick, petroleo}
  ]
    \node[caixa] (noticias) {Notícias\\InfoMoney\\(5.872 artigos)};
    \node[caixa, right=of noticias] (texto) {Representação textual\\\textit{embeddings} 1.024d\\ou sentimento 5d};
    \node[caixa, right=of texto] (merge) {\textit{Left join} por data\\+ \textit{forward-fill}};
    \node[caixa, right=of merge] (modelos) {Modelos\\BiLSTM, Transformer,\\XGBoost, TCN};
    \node[caixa, right=of modelos] (aval) {Avaliação\\AUC, \textit{split}\\cronológico 70/15/15};
    \node[caixa, below=of merge] (precos) {Yahoo Finance\\11 \textit{features} de preço};
    \draw[seta] (noticias) -- (texto);
    \draw[seta] (texto) -- (merge);
    \draw[seta] (merge) -- (modelos);
    \draw[seta] (modelos) -- (aval);
    \draw[seta] (precos) -- (merge);
  \end{tikzpicture}}
  \vspace{2ex}
  \begin{itemize}
    \item Dataset final: $\sim$1.200 dias úteis por ativo; janelas de 30 dias para modelos sequenciais
    \item \textit{Forward-fill}: propaga o passado, não vaza retorno futuro de mercado; resíduo intradiário \alert{não auditado}. A conclusão nula vale \alert{nas duas direções do viés}: se favorece o sentimento, é limite superior otimista; se for contra, o efeito real é ainda menor
  \end{itemize}
\end{frame}

% ---------------------------------------------------------------- 7. Etapa 3
\begin{frame}{Etapa 3 --- \textit{embeddings} genéricos: perto do acaso}
  \begin{columns}[T]
    \column{0.5\textwidth}
    \begin{itemize}
      \item Ollama \texttt{qwen3-embedding:4b}: 1.024 dimensões por artigo, média diária
      \item PCA $\to$ 32 componentes (61{,}4\% da variância), ajustado só no treino
      \item 1.227 dias $\times$ 1.035 \textit{features}
    \end{itemize}
    \vspace{1ex}
    \begin{block}{Leitura}
      Representação genérica de alta dimensionalidade carrega ruído
      semântico irrelevante ao domínio financeiro.
    \end{block}
    \column{0.46\textwidth}
    \begin{table}
      \footnotesize
      {\scriptsize\itshape ROC-AUC, janela única, $h=21$}\\[0.5ex]
      \begin{tabular}{lc}
        \toprule
        \textbf{Modelo} & \textbf{AUC} \\
        \midrule
        BiLSTM Original & 0{,}443 \\
        BiLSTM Reduzido & 0{,}505 \\
        Transformer     & 0{,}568 \\
        \textbf{XGBoost} & \textbf{0{,}610} \\
        \bottomrule
      \end{tabular}
    \end{table}
  \end{columns}
\end{frame}

% ---------------------------------------------------------------- 8. FinBERT
\begin{frame}{Etapa 4 --- sentimento específico: FinBERT-PT-BR}
  \begin{columns}[T]
    \column{0.5\textwidth}
    \textbf{O modelo}
    \begin{itemize}
      \item BERT com \textit{fine-tuning} em corpora financeiros em português (Souza et al., 2023)
      \item Calibrado para vocabulário do domínio: ``alavancagem'', ``\textit{guidance}'', ``\textit{rating}''
      \item 3 \textit{logits} por artigo: POSITIVO / NEGATIVO / NEUTRO
    \end{itemize}
    \column{0.46\textwidth}
    \textbf{5 \textit{features} diárias}
    \begin{itemize}
      \item \texttt{n\_articles}
      \item \texttt{mean\_logit\_pos / neg / neu}
      \item \texttt{mean\_sentiment}
    \end{itemize}
    \vspace{1ex}
    \begin{block}{Hipótese}
      Representação compacta e específica ao domínio é mais
      informativa do que 1.024 dimensões genéricas.
    \end{block}
  \end{columns}
\end{frame}

% ---------------------------------------------------------------- 9. AUC 0.709
\begin{frame}{O resultado que ``confirmou'' a hipótese}
  \begin{columns}[T]
    \column{0.5\textwidth}
    \begin{table}
      \footnotesize
      {\scriptsize\itshape Etapa 3 (\textit{embeddings}) vs Etapa 4 (FinBERT)}\\[0.5ex]
      \begin{tabular}{lccc}
        \toprule
        \textbf{Modelo} & \textbf{E3} & \textbf{E4} & $\boldsymbol{\Delta}$ \\
        \midrule
        BiLSTM Orig.    & 0{,}443 & 0{,}500 & $+$0{,}057 \\
        BiLSTM Red.     & 0{,}505 & 0{,}477 & $-$0{,}028 \\
        XGBoost         & 0{,}610 & 0{,}670 & $+$0{,}060 \\
        \textbf{Transformer} & 0{,}568 & \alert{\textbf{0{,}709}} & \textbf{$+$0{,}141} \\
        \bottomrule
      \end{tabular}
    \end{table}
    \column{0.46\textwidth}
    \begin{itemize}
      \item Transformer + FinBERT: \alert{ROC-AUC = 0{,}709}, acurácia 76{,}3\% \alert{($\approx$ classe majoritária da janela de TESTE 69{,}5\%; treino 59/41)}
      \item Narrativa coerente: ``sentimento específico supera \textit{embeddings} genéricos''
      \item Resultado compatível com a literatura internacional
    \end{itemize}
    \vspace{1ex}
    \centering\large
    Caso encerrado? \alert{\textbf{Não.}}
  \end{columns}
  \vspace{1ex}
  {\footnotesize Estimativa pontual, janela única, 1 semente, sem IC --- refutado como artefato nos slides 13--16.}
\end{frame}

% ---------------------------------------------------------------- 10. Alerta
\begin{frame}{Três sinais de alerta}
  \begin{enumerate}
    \item \textbf{Matriz de confusão assimétrica:} \textit{precision}(Desce) = 1{,}00 com \textit{recall}(Desce) = 0{,}20 --- previu ``Desce'' corretamente em só \alert{11 das 54 quedas reais}
    \item \textbf{Acurácia suspeita:} 76{,}3\% supera a classe majoritária da janela de teste (69{,}5\%; \textit{base rate} de treino 59/41) por só $\sim$7 pontos --- um ``sempre Sobe'' já acerta 69{,}5\%
    \item \textbf{Protocolo frágil:} estimativa pontual --- \alert{sem} intervalo de confiança, \alert{uma} semente, \alert{uma} janela temporal
  \end{enumerate}
  \vspace{1ex}
  \begin{block}{Decisão}
    Antes de aceitar o resultado: submetê-lo a protocolos progressivamente
    mais rigorosos. Se for real, sobrevive.
  \end{block}
\end{frame}

% ---------------------------------------------------------------- 11. Investigação
\begin{frame}{Investigação metodológica: 8 experimentos, 1.435 execuções}
  \begin{columns}[T]
    \column{0.55\textwidth}
    \begin{table}
      \footnotesize
      \begin{tabular}{lr}
        \toprule
        \textbf{Experimento} & \textbf{Runs} \\
        \midrule
        Baseline autoregressivo               & 4 \\
        Bootstrap CI                          & 1 \\
        Multi-seed (ITUB4)                    & 40 \\
        Multi-seed $\times$ multi-ticker      & 120 \\
        Ensemble + backtest                   & 20 \\
        Expanding-window CV                   & 145 \\
        VALE3 \textit{deep-dive}              & 880 \\
        Ablation PRICE/SENT/PRICE+SENT        & 225 \\
        \midrule
        \textbf{Total} & \textbf{1.435} \\
        \bottomrule
      \end{tabular}
    \end{table}
    \column{0.42\textwidth}
    \textbf{Protocolo (idêntico em todos)}
    \begin{itemize}
      \item Hiperparâmetros congelados; \texttt{StandardScaler} só no treino; bootstrap CI 1.000; Wilcoxon pareado
      \item Cada experimento = 1 pergunta pré-registrada; métrica nunca escolhida \textit{a posteriori}
      \item \alert{Oposto de \textit{p-hacking}:} as execuções \alert{destroem} o positivo, não selecionam o melhor. Sem correção para múltiplas comparações --- ela só tornaria o nulo mais forte
    \end{itemize}
  \end{columns}
  \vspace{1ex}
  {\footnotesize Centrais: \textbf{1.435} (8 pré-registrados). Exploratórios não geram $p$-valor.}
\end{frame}

% ---------------------------------------------------------------- 12. Baselines
\begin{frame}{Primeiro confronto: hierarquia de \textit{baselines}}
  \begin{table}
    \small
    \begin{tabular}{llc}
      \toprule
      \textbf{Preditor} & \textbf{Descrição} & \textbf{AUC [IC 95\%]} \\
      \midrule
      Classe majoritária    & Sempre prevê ``Sobe''           & 0{,}500 \\
      \textit{Coin flip}    & $P(\text{Sobe}) = \text{prior}$ & 0{,}500 \\
      Persistência ($h=21$) & Direção repete                  & 0{,}474 [0{,}405; 0{,}543] \\
      \textbf{Autoregressivo (XGB)} & \textbf{5 \textit{features} de preço (janela única)} & \textbf{0{,}658 [0{,}565; 0{,}744]} \\
      \bottomrule
    \end{tabular}
  \end{table}
  \vspace{1ex}
  \begin{itemize}
    \item XGBoost com só \texttt{return}, \texttt{lag\_1}, \texttt{lag\_5}, \texttt{Volume}, \texttt{std21} já atinge \textbf{0{,}658}
    \item A estimativa de janela única é tão incerta (IC largo) que \alert{não distingue} 0{,}709 do \textit{baseline} --- a separação só se resolve no multi-seed/CV (slides 13--16), onde a direção inverte
    \item \alert{AUC mede separabilidade, não retorno} --- e não reivindico nada sobre lucratividade (fora do escopo). O \textit{backtest} exploratório não deu sinal e foi descontinuado justamente por isso
  \end{itemize}
\end{frame}

% ---------------------------------------------------------------- 13. Multi-seed
\begin{frame}{Multi-seed: o colapso bimodal}
  \begin{columns}[T]
    \column{0.58\textwidth}
    \includegraphics[width=\textwidth]{multi_seed_histograms.png}\\[0.5ex]
    {\footnotesize\itshape Como ler: cada barra = 1 semente. \textcolor{verdeagua}{Azul} (meu modelo)
    espalhado de 0{,}08 a 0{,}93; \textcolor{coral}{vermelho} (\textit{baseline}) apertado em $\sim$0{,}65.
    A linha do 0{,}709 cai na \textbf{ponta direita} $\Rightarrow$ semente sortuda.}
    \column{0.38\textwidth}
    \begin{itemize}
      \item Semente 42: AUC cai de 0{,}709 para \alert{0{,}442} --- só trocando o sorteio
      \item 20 sementes: Transformer std = \alert{0{,}261} ($21\times$ o \textit{baseline}, que tem std 0{,}012)
      \item \textbf{Bimodal:} acumula nas pontas $\Rightarrow$ colapsa para ``sempre Sobe'' ou ``sempre Desce''
    \end{itemize}
  \end{columns}
\end{frame}

% ---------------------------------------------------------------- 14. Expanding CV
\begin{frame}{Expanding-window CV: a inversão}
  \begin{columns}[T]
    \column{0.55\textwidth}
    \includegraphics[width=\textwidth]{expanding_cv_overtime.png}\\[0.5ex]
    {\footnotesize\itshape Como ler: 3 ativos, eixo X = tempo (\textit{folds}).
    \textcolor{coral}{Vermelho} = preço, \textcolor{verdeagua}{azul} = meu modelo.
    Em ITUB4 e PETR4 o \textbf{vermelho fica por cima} $\Rightarrow$ o preço ganha.}
    \column{0.42\textwidth}
    {\footnotesize\alert{Leia primeiro: o ganho do sentimento na VALE3 nunca é estatisticamente significativo} (\textit{deep-dive} 880 \textit{runs}, $p=0{,}194$). Os $\Delta$ VALE3 abaixo vêm de \emph{experimentos distintos --- não comparáveis}.}\\[0.5ex]
    \begin{table}
      \footnotesize
      {\footnotesize\itshape AUC médio (5 \textit{folds} $\times$ 5 sementes)}\\[0.5ex]
      \begin{tabular}{lccc}
        \toprule
        \textbf{Ativo} & \textbf{XGB} & \textbf{Transf.} & $\boldsymbol{\Delta}$ \\
        \midrule
        ITUB4 & \textbf{0{,}700} & 0{,}445 & $-$0{,}255 \\
        PETR4 & \textbf{0{,}702} & 0{,}447 & $-$0{,}255 \\
        VALE3 & 0{,}599 & 0{,}635 & $+$0{,}036\,\alert{(n.s.; exp.\ distinto)} \\
        \midrule
        \textbf{Média} & \textbf{0{,}667} & 0{,}509 & $-$0{,}158 \\
        \bottomrule
      \end{tabular}
    \end{table}
    \begin{itemize}
      \item \textit{Baseline} 0{,}667 (média 5 \textit{folds} CV; cf.\ 0{,}658 da janela única, slide 12)
      \item Trocando 1 janela por 5 \textit{folds}: \alert{o resultado inverte}
      \item Diferença de 0{,}255 em ITUB4 $\approx 6\times$ o desvio do \textit{baseline}
    \end{itemize}
    \vspace{0.5ex}
    {\footnotesize \alert{VALE3: o ganho do sentimento nunca é estatisticamente significativo} (\textit{deep-dive} 880 \textit{runs}, $p=0{,}194$). Ganho aparente aqui $+$0{,}036; no \textit{deep-dive} (880 \textit{runs}) a VALE3 \alert{perde} (0{,}484 $<$ 0{,}544): magnitudes \emph{e até sinal} diferentes --- ver slide 15.}
  \end{columns}
\end{frame}

% ---------------------------------------------------------------- 15. VALE3
\begin{frame}{E a exceção VALE3? \textit{Deep-dive} com 880 execuções}
  \begin{columns}[T]
    \column{0.58\textwidth}
    \includegraphics[width=\textwidth]{vale3_deepdive_hist.png}\\[0.5ex]
    {\footnotesize\itshape Como ler: 880 execuções. \textcolor{verdeagua}{Azul} (meu modelo)
    acumula nas \textbf{duas pontas} (perto de 0 e de 1) = colapso. A média azul (0{,}484)
    fica à \textbf{esquerda} da \textcolor{coral}{vermelha} (0{,}544) $\Rightarrow$ perde até aqui.}
    \column{0.38\textwidth}
    {\footnotesize\alert{Leia primeiro: o ganho do sentimento na VALE3 nunca é estatisticamente significativo} ($p=0{,}194$, 880 \textit{runs}).}\\[0.5ex]
    \begin{itemize}
      \item Era o único ativo com ganho aparente ($+$0{,}036) na análise principal
      \item Reforço: \textbf{52 \textit{folds} $\times$ 10 sementes} = 880 execuções
      \item Resultado: \alert{não significativo} (Wilcoxon $p = 0{,}194$) --- a exceção morreu
    \end{itemize}
    \vspace{0.5ex}
    {\footnotesize \alert{VALE3: o ganho do sentimento nunca é estatisticamente significativo} (\textit{deep-dive} 880 \textit{runs}, $p=0{,}194$). Três experimentos distintos (modelos/contrastes/$n$ diferentes nos slides 14/15/16) $\Rightarrow$ magnitudes diferentes; a instabilidade do sinal é a prova de que não é sinal.}
  \end{columns}
\end{frame}

% ---------------------------------------------------------------- 16. Ablation
\begin{frame}{Ablation: o sentimento adiciona algo ao preço?}
  \begin{columns}[T]
    \column{0.52\textwidth}
    \includegraphics[width=\textwidth]{ablation_boxplot.png}\\[0.5ex]
    {\footnotesize\itshape Como ler: 3 caixas por ativo. \textcolor{verdeagua}{Verde} (só sentimento)
    fica \textbf{na linha do acaso ou abaixo}. \textcolor{coral}{Vermelho} (preço) $\approx$ azul (preço+sent.)
    $\Rightarrow$ somar sentimento quase não move nada.}
    \column{0.45\textwidth}
    {\footnotesize\alert{Leia primeiro: o ganho do sentimento na VALE3 nunca é estatisticamente significativo} ($p=0{,}194$, 880 \textit{runs}). O $+$0{,}058 abaixo é \emph{experimento distinto --- não comparável}.}\\[0.5ex]
    \begin{table}
      \footnotesize
      {\footnotesize\itshape XGBoost, 225 treinamentos, $h=21$}\\[0.5ex]
      \resizebox{\linewidth}{!}{%
      \begin{tabular}{lcccc}
        \toprule
        \textbf{Ativo} & \textbf{PRICE} & \textbf{SENT} & \textbf{P+S} & $\boldsymbol{\Delta}$ \\
        \midrule
        ITUB4 & 0{,}684 & 0{,}436 & 0{,}651 & $-$0{,}033 \\
        PETR4 & 0{,}692 & 0{,}494 & 0{,}676 & $-$0{,}016 \\
        VALE3 & 0{,}609 & 0{,}510 & 0{,}667 & $+$0{,}058\,\alert{(n.s.; exp.\ distinto)} \\
        \midrule
        \textbf{Média} & \textbf{0{,}662} & 0{,}480 & 0{,}665 & \alert{\textbf{$+$0{,}003}} \\
        \bottomrule
      \end{tabular}}
    \end{table}
    \begin{center}\large\alert{$\Delta = +0{,}003$}\ \ ($p = 0{,}49$)\end{center}
    \begin{itemize}
      \item IC 95\% $[-0{,}012;\ +0{,}018]$ cruza o zero; SENT sozinho \alert{abaixo do acaso}
      \item \emph{Não detectável neste $n$}, não ``efeito zero'': MDE 0{,}13--0{,}22 em AUC (sentimento 44--74$\times$ abaixo) --- ver slide 17
    \end{itemize}
    \vspace{0.5ex}
    {\footnotesize \alert{VALE3: o ganho do sentimento nunca é estatisticamente significativo.} O $+$0{,}058 desta ablation é isolado; significância formal só no \textit{deep-dive} de 880 \textit{runs} ($p=0{,}194$).}
  \end{columns}
\end{frame}

% ---------------------------------------------------------------- 17. Síntese
\begin{frame}{Síntese da investigação}
  \begin{block}{Conclusão unificada}
    Sob avaliação correta (multi-\textit{fold}, multi-semente, IC, \textit{ablation}),
    as \textit{features} de sentimento FinBERT-PT-BR \alert{não adicionam sinal
    preditivo mensurável} ao \textit{baseline} autoregressivo --- em nenhum
    dos três ativos.
  \end{block}
  \vspace{1ex}
  \textbf{O AUC de 0{,}709 era um artefato de janela única, produzido por:}
  \begin{enumerate}
    \item Variância de semente (std 0{,}261 --- colapso bimodal do Transformer)
    \item Viés de amostragem da janela de teste específica
    \item Ausência de intervalos de confiança que revelariam a incerteza
  \end{enumerate}
  \vspace{1ex}
  \footnotesize Escopo: conclusão vale para a representação adotada (5 \textit{features} diárias FinBERT-PT-BR),
  não para toda forma de incorporar texto. Resultado é ``nenhum efeito detectável neste tamanho amostral''.
\end{frame}

% ---------------------------------------------------------------- 18. Contribuições
\begin{frame}{Contribuições e protocolos mínimos propostos}
  \begin{columns}[T]
    \column{0.46\textwidth}
    \textbf{Três contribuições}
    \begin{enumerate}
      \item Pipeline replicável e versionado (InfoMoney + FinBERT-PT-BR + Yahoo Finance)
      \item Primeira demonstração quantificada do viés de janela única em dados brasileiros (1.435 execuções)
      \item Proposição de protocolos mínimos para ML financeiro
    \end{enumerate}
    \column{0.5\textwidth}
    \textbf{Seis práticas mínimas}
    \begin{enumerate}
      \item Reportar IC \textit{bootstrap}
      \item Treinar com $\geq$10 sementes
      \item CV \textit{expanding-window}
      \item Comparar contra \textit{baseline} autoregressivo
      \item Monitorar distribuição de predições
      \item Auditar correlação validação-teste
    \end{enumerate}
  \end{columns}
  \vspace{1ex}
  \footnotesize Práticas alinhadas a López de Prado (2018); a contribuição é mostrar,
  em caso concreto, a magnitude do dano de ignorá-las.
\end{frame}

% ---------------------------------------------------------------- 19. Limitações
\begin{frame}{Limitações e trabalhos futuros}
  \begin{columns}[T]
    \column{0.48\textwidth}
    \textbf{Limitações}
    \begin{itemize}
      \item Fonte única (InfoMoney, viés varejo)
      \item 3 ativos de grande capitalização
      \item Horizontes restritos ($h=5$, $h=21$)
      \item Representação textual específica (5 \textit{features} agregadas)
      \item \textit{Look-ahead} residual do \textit{forward-fill}: propaga passado (não vaza futuro de mercado). Resíduo é apenas intradiário (notícia pós-17h atribuída ao dia $t$), escopo limitado e \alert{não auditado}. A conclusão nula é conservadora sob o cenário adverso e apenas reforçada sob o favorável --- em nenhuma das direções a conclusão muda
      \item Instabilidade arquitetural do Transformer
    \end{itemize}
    \column{0.48\textwidth}
    \textbf{Trabalhos futuros}
    \begin{itemize}
      \item Diversificar fontes (CVM, \textit{analyst reports}, redes sociais)
      \item Arquiteturas menos propensas a colapso bimodal
      \item Outros mercados emergentes
      \item Estudo formal da anticorrelação validação-teste
      \item Métricas robustas a desbalanceamento (\textit{balanced accuracy}, MCC)
    \end{itemize}
  \end{columns}
\end{frame}

% ---------------------------------------------------------------- 20. Fechamento
\begin{frame}{Mensagem final}
  \begin{block}{}
    \centering\large
    A contribuição central não é descobrir que o sentimento não funciona ---\\
    é demonstrar \alert{quão fácil é produzir evidência aparente de que funciona},\\
    usando protocolos que a comunidade considera aceitáveis.
  \end{block}
  \vspace{2ex}
  \begin{itemize}
    \item Viés de confirmação opera de forma invisível: o 0{,}709 ``fazia sentido''
    \item Janela única + semente única + ausência de IC: receita para falso positivo
    \item Autocorreção documentada $>$ resultado positivo frágil
  \end{itemize}
  \vspace{2ex}
  \centering
  \textbf{Obrigado!}\\[1ex]
  \footnotesize André Takeo Loschner Fujiwara --- Ciência de Dados e IA, PUC-SP\\
  Código e dados: repositório versionado do projeto
\end{frame}

\end{document}
